\subsection{弗罗贝尼乌斯--舒尔指标与拉森择一定理}

回忆（LIE, VIII, 第 7 节第 6 号定义 2）：若 G 在有限维复向量空间 E 中的不可约表示存在一个在 G 下不变的非零对称双线性形式（分别地，存在一个在 G 下不变的非零交错双线性形式），则称该表示为正交型（分别为辛型）。若 E 上不存在在 G 下不变的非零双线性形式，则称其为复型。

当 E 中的不可约表示为正交型（分别为辛型）时，E 上 G-不变对称（分别为交错）双线性形式的空间是一维的，并且每个非零 G-不变双线性形式都是分离的。

正交型表示有时称为实型表示，辛型表示有时称为四元数型表示（参见 LIE, IX, 附录 II）。

令 \(\pi\) 为 G 在复向量空间 E 中的不可约表示。上述三种情形由量

\[
\mathrm{FS} (\pi) = \int_ {\mathrm{G}} \chi_ {\pi} (g ^ {2}) d \mu (g),
\]

区分；称其为 π 的弗罗贝尼乌斯--舒尔指标。事实上有

\[
\operatorname{FS} (\pi) = \left\{ \begin{array}{l l} 1 & \text{若 }\pi\text{ 属于正交型，} \\ - 1 & \text{若 }\pi\text{ 属于辛型，} \\ 0 & \text{若 }\pi\text{ 属于复型。} \end{array} \right.
\]

（LIE, IX, 附录 2，第 103 页命题 3 及第 105 页命题 4）。

当 G 为李群时，可以利用 LIE, IX, 第 69 页命题 1 计算 \(\mathrm{FS}(\pi)\)。

定义 3. --- 令 \(\varrho\) 为 G 在希尔伯特空间 E 中的有限维酉表示。令 \(k\) 为正整数。称表示 \(\overline{\varrho}^{\otimes k} \otimes \varrho^{\otimes k}\) 的 G-不变元素子空间的维数为 \(2k\) 阶的 \(\varrho\) 的绝对矩，记为 \(\mathrm{M}_{2k}(\varrho)\)。

定理 3（拉森择一定理）。--- 设 G 是无限群。令 π 为 G 在维数 ≥ 2 的有限维希尔伯特空间 E 中的忠实酉表示。

\begin{enumerate}
\def\labelenumi{\alph{enumi})}
\item
  假设 G 的导出群是无限的。有 \(M_{4}(\pi) \geqslant 2\)，且等号成立当且仅当 G 包含 SU(E)；
\item
  假设 \(\dim(E) \geqslant 3\)，\(\pi\) 为正交型或辛型，且 \(\dim(E) \neq 4\)（在 \(\pi\) 为正交型时）。则 \(M_4(\pi) \geqslant 3\)，且等号成立当且仅当 G 的复化李代数等于 E 上某个分离的 G-不变双线性形式的李代数。*当且仅当 G 的单位分支 \(G_0\) 是这种双线性形式的自同构群的极大紧子群时，情形才成立。*
\end{enumerate}

证明中要用到下列引理。

引理 5. --- 令 \(\pi\) 为 G 在有限维希尔伯特空间 E 中的酉表示。令 \((\varrho_{i})_{i\in\mathrm{I}}\) 为 G 的非零酉表示族，\((n_{i})_{i\in\mathrm{I}}\) 为满足 \(\geqslant 1\) 的整数族，使得 G 的表示 \(\overline{\pi}\otimes\pi\)（分别地，表示 \(\pi\otimes\pi\)）同构于直和 \(\bigoplus_{i\in\mathrm{I}}\varrho_{i}^{n_{i}}\)。则有

\[
\mathrm{M} _ {4} (\pi) \geqslant \sum_ {i \in \mathrm{I}} n _ {i} ^ {2}
\]

等号成立当且仅当表示 \(\varrho_{i}\) 都不可约且两两不同构。

记 \(\chi_{i}\) 为 \(\varrho_{i}\) 的特征标，其中 \(i\in I\)。有

\[
| \chi_ {\pi} | ^ {2} = \chi_ {\overline {{\pi}} \otimes \pi} = \sum_ {i \in \mathrm{I}} n _ {i} \chi_ {i}, \quad (\text { resp. } \chi_ {\pi} ^ {2} = \chi_ {\pi \otimes \pi} = \sum_ {i \in \mathrm{I}} n _ {i} \chi_ {i}).
\]

根据定义和 V, 第 466 页推论 3，有

\[
\mathrm{M} _ {4} (\pi) = \int_ {\mathrm{G}} | \chi_ {\pi} | ^ {4} d \mu ,
\]

从而在两种情形下均有公式

\[
\mathrm{M} _ {4} (\pi) = \sum_ {i, j} n _ {i} n _ {j} \int_ {\mathrm{G}} \chi_ {i} \overline {{\chi}} _ {j} d \mu = \sum_ {i, j} n _ {i} n _ {j} \dim \mathrm{Hom} _ {\mathrm{G}} (\varrho_ {i}, \varrho_ {j})
\]

（V, 第 459 页推论 4, b)）。令 \(i \in I\)。由于表示 \(\varrho_{i}\) 非零，空间 \(\mathrm{Hom}_{\mathrm{G}}(\varrho_{i}, \varrho_{i})\) 非零，于是得到下界

\[
\mathrm{M} _ {4} (\pi) \geqslant \sum_ {i \in \mathrm{I}} n _ {i} ^ {2}.
\]

此外，由于 \(n_{i} \geqslant 1\)，等号成立当且仅当 \(\dim \operatorname{Hom}_{\mathrm{G}}(\varrho_{i}, \varrho_{j}) = 0\)（当 \(i \neq j\) 时），并且对每个 i 都有 \(\dim \operatorname{Hom}_{\mathrm{G}}(\varrho_{i}, \varrho_{i}) = 1\)。第二个条件成立当且仅当表示 \(\varrho_{i}\) 不可约（上引处）。此时根据舒尔引理（V, 第 387 页推论 2），第一个条件成立当且仅当表示 \(\varrho_{i}\) 两两不同构。

若 E 是交换环 A 上的模，则称 E 的对称平方（分别为外平方）为 A-模 \(S^{2}(E)\)（分别为 \(\wedge^{2}E\)）。

引理 6. --- 令 \(q\) 为维数 \(\geqslant 3\) 的有限维复向量空间 E 上的分离双线性形式。

\begin{enumerate}
\def\labelenumi{\alph{enumi})}
\item
  若 \(q\) 是对称的，则正交李代数 \(\mathfrak{so}(q)\) 的伴随表示同构于 \(\wedge^2\mathrm{E}\)，即自然表示 \(\mathfrak{so}(q) \to \mathfrak{gl}(\mathrm{E})\) 的外平方；
\item
  若 q 是交错的，则辛李代数 \(\mathfrak{sp}(q)\) 的伴随表示同构于 \(\mathsf{S}^2 (\mathrm{E})\)，即自然表示 \(\mathfrak{sp}(q)\to \mathfrak{gl}(\mathrm{E})\) 的对称平方。
\end{enumerate}

在 C 上赋予恒等对合自同构。于是双线性形式 q 是 \(\varepsilon\)-埃尔米特的（A, IX, 第 3 节第 1 号定义 1），其中情形 a) 中 \(\varepsilon = 1\)，情形 b) 中 \(\varepsilon = -1\)。对于每个 \(u \in \text{End}(\text{E})\)，记 \({}^{t}u\) 为关于 q 的 u 的伴随。于是有 \(q(x, u(y)) = q(^{t}u(x), y)\)，对所有 \((x, y) \in \text{E} \times \text{E}\) 成立。令 v 为 E \(\otimes\) E 的唯一自同构，使得 \(v(x \otimes y) = y \otimes x\) 对所有 \((x, y) \in \text{E}^{2}\) 成立。

存在唯一的线性映射 \(w\) 从 \(\mathrm{E} \otimes \mathrm{E}\) 到 \(\operatorname{End}(\mathrm{E})\)，使得 \(w(a \otimes b)\) 是由 \(x \mapsto q(a, x)b\) 定义的线性映射，对所有 \((a, b, x) \in \mathrm{E}^3\) 成立。映射 \(w\) 是同构。对于每个 \(s \in \mathrm{E} \otimes \mathrm{E}\)，有

\[
{ } ^ { t } ( w ( s ) ) = \varepsilon w ( v ( s ) ) .\tag{3}
\]

事实上，对于 \((x,y)\) 和 \((a,b)\) 在 \(\mathrm{E} \times \mathrm{E}\) 中的所有情形，可得

\[
q (x, w (a \otimes b) y) = q (a, y) q (x, b) = \varepsilon q (b, x) q (a, y) = q (\varepsilon w (b \otimes a) x, y).
\]

最后，记 \(x \mapsto x^*\) 为由 q 导出的从 E 到 \(\mathrm{E}'\) 的同构，即由 \(q\) 给出，满足 \(\langle x^*, y \rangle = q(x, y)\)，对所有 \((x, y) \in \mathrm{E}^2\) 成立。

记号确定后，令 H 为 q 的正交群（分别为辛群）。它是 \(\mathbf{GL}(\mathrm{E})\) 的李子群，其李代数 h 是由 \(u \in \operatorname{End}(\operatorname{E})\) 满足 \(^{t}u = -u\) 的 End(E) 子空间（LIE, III, 第 146 页推论 1）。由于形式 q 在 H 下不变，故有

\[
\langle (g a) ^ {*}, b \rangle = q (g a, b) = \langle a, g ^ {- 1} b \rangle
\]

对所有 \(g \in H\) 及每个 \((a, b) \in E \times E\) 都成立。

在 End(E) 上赋予 H 的伴随表示 Ad。线性映射 \(w\) 是 H 的表示之间的态射：事实上，对每个 \(g \in \mathrm{H}\) 及每个 \((a, b, x) \in \mathrm{E}^3\)，有

\[
\begin{array}{c} w (g   (a \otimes b)) (x) = \langle (g a) ^ {*}, x \rangle   g b = \langle a, g ^ {- 1} x \rangle   g b \\ = g   (\langle a, g ^ {- 1} x \rangle   b) = \mathrm{Ad} (g) (w (a \otimes b)) x, \end{array}
\]

由线性性即得结论。

由于 v 也是 H 的表示之间的态射，线性映射 \(\theta = w - \varepsilon w \circ v\) 也具有这一性质；因此它是 h-模之间的态射。

记 \(F \subset E \otimes E\) 为由元素 \(s \in E \otimes E\) 且满足 \(v(s) = -\varepsilon s\) 所成的子空间。若 \(\varepsilon = -1\)，则从 \(E \otimes E\) 到 E 的对称平方的典范映射在 F 上的限制是 C-向量空间同构（A, III, 第 72 页）；若 \(\varepsilon = 1\)，则从 \(E \otimes E\) 到外平方的典范映射在 F 上的限制是 C-向量空间同构（上引处，第 82 页）。

F 在 \(\theta\) 下的像包含于 \(\mathfrak{h}\) 中：事实上，对每个 \(s\in\mathrm{F}\)，公式 (3) 蕴含

\[
{ } ^ { t } \theta ( s ) = { } ^ { t } w ( s ) - \varepsilon { } ^ { t } w ( v ( s ) ) = \varepsilon w ( v ( s ) ) - w ( s ) = - \theta ( s ) .
\]

按 F 的定义，映射 \(\theta\) 在 F 上的限制与 \(2w\) 的限制重合，因而线性映射 \(\theta\) 在 F 上是单射。由于 \(\dim(F) = \dim(h)\)。由（A, III, 第 75 页定理 1 和第 87 页推论 1，以及 LIE, VIII, 第 192 页和第 201 页），可知 \(\theta\) 通过取子空间诱导出从 F 到 \(h\) 的 \(h\)-模同构，故引理得证。

引理 7. --- 令 \(H_{2}\) 为实李群，令 \(H_{1}\) 为 \(H_{2}\) 的闭子群。\(H_{1}\) 的复化李代数可识别为 \(H_{1}\) 在 \(H_{2}\) 的复化李代数上的伴随表示的一个子表示。

事实上，限制到 \(H_{1}\) 的 \(H_{2}\) 李代数上的伴随表示，正可识别为 \(H_{1}\) 的伴随表示。

引理 8. --- 令 E 为维数 n 的复希尔伯特空间。

\begin{enumerate}
\def\labelenumi{\alph{enumi})}
\item
  群 \(\mathbf{SU}(\mathrm{E})\) 和 \(\mathbf{U}(\mathrm{E})\) 是连通的；
\item
  SU(E) 的导出群等于 SU(E)；
\item
  U(E) 的导出子群等于 SU(E)。
\end{enumerate}

不妨设 \(n \geqslant 1\) 且 \(E = C^{n}\)。

令 A 为 \(\mathbf{U}(\mathrm{E})\) 中由对角矩阵组成的子群；它同胚于 \(U^{n}\)，因而是连通的（TG, VI, 第 11 页推论 2 以及 I, 第 83 页命题 8）。它与 \(\mathbf{SU}(\mathrm{E})\) 的交同胚于 \(U^{n-1}\)，所以也是连通的。

由 IV, 第 149 页定理 1，群 \(\mathbf{U}(\mathrm{E})\)（分别为 \(\mathbf{SU}(\mathrm{E})\)）是连通子集 \(g\mathrm{Ag}^{-1}\)（\(g \in \mathrm{G}\)）的并（分别是连通子集 \(g(\mathrm{A} \cap \mathbf{SU}(\mathrm{E}))g^{-1}\) 的并）；由于单位元属于这些集合中的每一个，空间 \(\mathbf{U}(\mathrm{E})\)（分别为 \(\mathbf{SU}(\mathrm{E})\)）是连通的（TG, I, 第 81 页命题 2）。

证明 \(b\)。当 \(n = 1\) 时断言成立，因为此时 \(\mathbf{SU}(\mathrm{E})\) 退化为单位元。因此设 \(n \geqslant 2\)。此时 \(\mathbf{SU}(\mathrm{E})\) 的李代数是单的（LIE, IX, 第 20 页，第 3 节第 4 号），所以 \(\mathbf{SU}(\mathrm{E})\) 的导出群在 \(\mathbf{SU}(\mathrm{E})\) 中具有有限指数。由于 \(\mathbf{SU}(\mathrm{E})\) 由 \(a\) 可知连通，故得到所需结果。

断言 \(c\) 由 \(b\) 得出，因为 \(\mathbf{U}(\mathrm{E})\) 的导出群包含于 \(\mathbf{SU}(\mathrm{E})\)。

现在证明定理 3。记 n 为希尔伯特空间 E 的维数。由于表示 \(\pi\) 忠实，不妨设 G 是 U(E) 的紧子群。群 G 在实李群 U(E) 中是闭的，因而是紧实李群（LIE, III, 第 8 节第 2 号定理 2）。根据假设，G 是无限的，故其维数 \(\geqslant 1\)。令 g 为其李代数；它非零。

证明 \(a\)。假设导出群 D(G) 是无限的。有 G-不变分解 \(\operatorname{End}(\mathrm{E}) = \mathbf{C}1_{\mathrm{E}} \oplus \mathfrak{sl}(\mathrm{E})\)。G 在 \(\mathbf{C}1_{\mathrm{E}}\) 和 \(\mathfrak{sl}(\mathrm{E})\) 上的表示不同构，因为 \(\dim \mathfrak{sl}(\mathrm{E}) \geqslant 2\)。由引理 5，\(\mathrm{M}_4(\pi) \geqslant 2\)，且等号成立当且仅当 G 在 \(\mathfrak{sl}(\mathrm{E})\) 中的表示不可约。现在，G 的导出群包含于 \(\mathbf{SL}(\mathrm{E})\)，所以 D(G) 的李代数的复化可识别为 \(\mathfrak{sl}(\mathrm{E})\) 的一个子空间，而后者是 G 在 \(\mathfrak{sl}(\mathrm{E})\) 上的表示的一个子表示（引理 7）。由于 D(G) 无限，该子表示非零。因此，\(\mathrm{M}_4(\pi) = 2\) 当且仅当 \((\mathcal{D}\mathfrak{g})_{(\mathbf{C})} = \mathfrak{sl}(\mathrm{E})\)。由于 D(G) 包含于 \(\mathbf{SU}(\mathrm{E})\)，且 \(\mathbf{SU}(\mathrm{E}) = \mathrm{D}(\mathbf{SU}(\mathrm{E}))\)（引理 8），这个条件等价于 \(\mathbf{SU}(\mathrm{E}) \subset \mathbf{G}\)。

为证明断言 \(b\)，沿用 LIE, VIII, 第 13 节第 2--4 号的记号。

设 \(\pi\) 为辛型且 \(\dim(\mathrm{E}) \geqslant 3\)。令 \(q\) 为 E 上非零的 G-不变交错双线性形式；它是分离的（LIE, IX, 第 103 页附录 2 命题 3），且 G 是辛群 \(\mathbf{Sp}(q)\) 的紧子群。记 \(q^{*} \in \wedge^{2}\mathrm{E}\) 为交错形式 \(q\) 所对应的元素。G 的复化李代数 \(\mathfrak{g}_{(\mathbf{C})}\) 包含于 \(\mathfrak{sp}(q)\)。现在，由于 \(\dim(\mathrm{E}) \geqslant 3\)，\(\mathfrak{sp}(q)\) 在 \(\mathrm{E} \otimes \mathrm{E}\) 中的表示具有如下直和分解

\[
\mathrm{E} \otimes \mathrm{E} = \mathsf {S} ^ {2} (\mathrm{E}) \oplus \wedge^ {2} \mathrm{E} = \mathsf {S} ^ {2} (\mathrm{E}) \oplus \mathrm{E} _ {2} \oplus \mathbf {C} q ^ {*},
\]

其中 \(E_{2}\) 是 \(\mathfrak{sp}(q)\) 的第二基本表示（LIE, VIII, 第 202 页，第 13 节第 3 号 (IV)）。表示 \(E_{2}\) 和 \(C q^{*}\) 是 \(\mathfrak{sp}(q)\) 的不可约表示（上引处），而表示 \(S^{2}(E)\) 非零。

因此由引理 5，\(\mathrm{M}_{4}(\pi) \geqslant 3\)，且等号成立当且仅当 G 在 \(\mathsf{S}^{2}(\mathrm{E})\)、\(E_{2}\) 和 \(Cq^{*}\) 中的表示都不可约且两两不同构。

设 \(M_{4}(\pi)=3\)。由于表示 \(S^{2}(E)\) 可识别为 \(\mathfrak{sp}(q)\) 的伴随表示（引理 6），它包含 G 的伴随表示的复化作为子表示（引理 7）。因此有 \(\mathfrak{g}_{(\mathbf{C})}=\mathfrak{sp}(q)\)。

反之，设 \(\mathfrak{g}_{(\mathbf{C})} = \mathfrak{sp}(q)\)。\(\mathfrak{sp}(q)\) 的伴随表示和表示 \(\mathrm{E}_2\)，后者属于 \(\mathfrak{sp}(q)\)，都不可约，其维数分别为 \(n(n + 1)/2\) 和 \(n(n - 1)/2 - 1\)（LIE, VIII, 第 202 页，第 13 节第 3 号 (IV)）；这两个维数不同且均 \(\geqslant 2\)，因为 \(n \geqslant 3\)，从而 \(\mathrm{M}_4(\pi) = 3\)。

最后，设 \(\pi\) 为实型，且 \(\dim(\mathrm{E}) \geqslant 3\) 以及 \(\dim(\mathrm{E}) \neq 4\)。令 \(q\) 为 \(\mathrm{E}\) 上的非零对称双线性形式，且在 G 下不变；它是分离的（LIE, IX, 第 103 页附录 2 命题 3），并且 \(\mathrm{G}\) 是正交群 \(\mathbf{O}(q)\) 的紧子群。记 \(q^{*} \in \mathbb{S}^{2}(\mathrm{E})\) 为与 \(q\) 对应的元素。

G 的复化李代数 \(\mathfrak{g}_{(\mathbf{C})}\) 包含于 \(\mathfrak{so}(q)\)。\(\mathfrak{so}(q)\) 在 E \(\otimes\) E 中的表示具有如下直和分解

\[
\mathrm{E} \otimes \mathrm{E} = \wedge^ {2} \mathrm{E} \oplus \mathsf {S} ^ {2} (\mathrm{E}) = \wedge^ {2} \mathrm{E} \oplus \mathsf {S} _ {0} ^ {2} (\mathrm{E}) \oplus \mathbf {C} q ^ {*},
\]

其中 \(\mathsf{S}_{0}^{2}(\mathbf{E})\) 是 \(q^{*}\) 在 \(\mathsf{S}^{2}(\mathbf{E})\) 中的正交补。这些表示的维数至少为 2，因为 \(\dim(\mathbf{E})\geqslant3\)。由引理 5，\(\mathrm{M}_{4}(\pi)\geqslant3\)，且等号成立当且仅当 G 在 \(\wedge^{2}\mathrm{E}\) 和 \(\mathsf{S}_{0}^{2}(\mathbf{E})\) 中的表示不可约且不同构。

设 \(\mathrm{M}_{4}(\pi)=3\)。由于表示 \(\wedge^{2}E\) 可识别为 \(\mathfrak{so}(q)\) 的伴随表示（引理 6），它包含 G 的伴随表示的复化作为子表示（引理 7）。因此条件 \(\mathrm{M}_{4}(\pi)=3\) 蕴含 \(\mathfrak{g}_{(\mathbf{C})}=\mathfrak{so}(q)\)。

反之，设 \(\mathfrak{g}_{(\mathbf{C})} = \mathfrak{so}(q)\)。\(\mathfrak{so}(q)\) 的伴随表示是不可约的，因为 \(\dim(\mathrm{E}) \neq 4\)（LIE, VIII, 第 13 节，第 193 页 (I) 及第 206 页 (I)），且维数为 \(n(n-1)/2\)。表示 \(\mathsf{S}_0^2(\mathrm{E})\) 是 \(\mathfrak{so}(q)\) 的表示，其最高权为 \(2\varpi_1\)。将其维数与由 Weyl 公式计算的 \(\mathfrak{so}(q)\) 的最高权为 \(2\varpi_1\) 的不可约表示的维数比较（参见 LIE, VIII, 第 9 节第 2 号定理 2 以及 LIE, VI, 图版 II 和 IV），可验证 \(\mathsf{S}_0^2(\mathrm{E})\) 是不可约的。因此得到 \(\mathrm{M}_4(\pi) = 3\)，当 \(\mathfrak{g}_{(\mathbf{C})} = \mathfrak{so}(q)\) 时。

注记。--- 1) “G 是无限的”这一条件对于定理 3 是必要的（V, 第 507 页习题 9）。

\begin{enumerate}
\def\labelenumi{\arabic{enumi})}
\setcounter{enumi}{1}
\item
  令 \(n \in \mathbf{N}\)。对于紧群的每个酉表示 \(\varrho\)，若它作用在维数为 \(n\) 的希尔伯特空间中，则有 \(\mathrm{M}_{4}(\varrho) \leqslant n^{2}\)；等号可能成立（V, 第 508 页习题 15）。
\item
  可以证明（参见 R. GURALNICK 和 P.H. TIEP，《Decomposition of small tensor powers and Larsen's conjecture》，Representation Theory 9 (2005), 138--208）：若假设 E 的维数 \(\geqslant 7\)，并将假设 \(M_{4}(\pi) = 2\)（分别为 \(M_{4}(\pi) = 3\)）替换为 \(M_{8}(\pi) = 24\)（分别为 \(M_{8}(\pi) = 105\)），则可以省略 G 为无限这一条件。
\end{enumerate}

